% Appendix C: Numerical evolution of the linearised system (~2000 w).
% Per genre_conventions_software.md §(j). Renamed from "Numerical stage"
% — the title now describes what the appendix actually does (evolves the
% linearised system numerically), matching the BHL convention of
% content-bearing names. The Fourier modal solver receives the longest
% treatment because it is the primary backend; the others are documented
% briefly.

\section{Numerical evolution of the linearised system}%
\label{NumericalEvolution}

\paragraph*{Solver selection hierarchy}
% TODO: opener — modal --> IDA --> CVODE --> leapfrog --> scipy;
% auto-selection logic: which equation properties trigger each backend.

\subsection{Fourier modal solver}%
\label{ModalSolver}
% Target ~800 w — primary backend, longest treatment.

\paragraph*{Eigendecomposition}
% TODO: constant-coefficient case — exact matrix exponential. State what
% "machine precision" means in this context.

\paragraph*{Position-dependent coefficients}
% TODO: Krylov \texttt{expm\_multiply}.

\paragraph*{Constraint elimination}
% TODO: Fourier Schur complement.

\paragraph*{Mathematical foundations}
% TODO: cite \cite{higham2008functions} and \cite{molervanloan2003nineteen}
% for matrix-exponential methods.

% Figure C.1 — Fourier modal solver convergence. Referenced from App D.
\begin{figure}[t]
  \centering
  % TODO: \includegraphics{figC1_solver_convergence}
  % TODO caption: "Fourier modal solver convergence: relative error vs grid
  % resolution for Einstein–Maxwell at $\Bzero = [\text{value}]$,
  % $D = [\text{value}]$. The solver attains machine precision
  % ($\sim 10^{-14}$) at modest resolution, demonstrating the
  % spectral-precision behaviour expected for a constant-coefficient
  % system. This figure also satisfies the App D convergence requirement."
  \caption{TODO: solver convergence vs resolution.}
  \label{SolverConvergence}
\end{figure}

\subsection{Analytical Jacobian}%
\label{AnalyticalJacobian}
% Target ~400 w.

\paragraph*{Three-tier structure}
% TODO: dense (N <= 2K), sparse (2K < N <= 200K), GMRES (N > 200K). One
% sentence on when each is active.

\paragraph*{Speedup factors}
% TODO: quote the measured speedups (5.3x dense; 2.5x IDA sparse;
% 1.2–1.4x CVODE sparse) from solver_optimizations.tex.

\subsection{Other backends}%
\label{OtherBackends}
% Target ~200 w total.

\paragraph*{IDA}
% TODO: 3–4 sentences on DAE / constraint handling.

\paragraph*{CVODE}
% TODO: 3–4 sentences on adaptive ODE for waves.

\paragraph*{Leapfrog}
% TODO: 3–4 sentences on symplectic integration.

\paragraph*{scipy}
% TODO: 3–4 sentences on general-purpose fallback.

\paragraph*{JSON-to-solver interface}
% TODO: closing paragraph with two distinct moves —
% (i) coefficient-matrix EXTRACTION from the JSON spec (which fields are
% read, how they are interpreted); (ii) solver DISPATCH (how
% auto-selection decides which backend to call given the extracted
% properties). One sentence each.
